\chapter{Chapter09 — Learning Consequences}
%%% generated by the blueprint pipeline from FABL.Chapter09.LearningConsequences
\section{Overview}

Declarations of the Lean module \texttt{FABL.Chapter09.LearningConsequences}.

\section{Declarations}

\begin{definition}[Positive rational influence bound]
\lean{FABL.PositiveInfluenceBound}
\statementsource{lmn-constant-depth-fourier}{
  pdf-63759bd909d5-p011-b006,
  pdf-63759bd909d5-p011-b007,
  pdf-63759bd909d5-p011-b009
}
\statementsource{Analysis-of-Boolean-Functions-by-Ryan-ODonnell}{
  pdf-ffbe9c7b3306-p069-b004,
  pdf-ffbe9c7b3306-p265-b004
}
\label{FABL.PositiveInfluenceBound}
\leanok
A positive influence bound is a rational number $q$ satisfying $q>0$.
\end{definition}

\begin{definition}[Internal learning error]
\lean{FABL.learningConsequencesInternalError}
\label{FABL.learningConsequencesInternalError}
\leanok
\uses{FABL.PositiveLearningParameter}
For a rational learning parameter $\varepsilon$ with $0<\varepsilon\le 1/2$, define the
internal error to be $\varepsilon/8$.
\end{definition}

\begin{theorem}[Positivity of the internal error]
\lean{FABL.learningConsequencesInternalError_pos}
\label{FABL.learningConsequencesInternalError_pos}
\leanok
\uses{FABL.PositiveLearningParameter, FABL.learningConsequencesInternalError}
\difficulty{1}
For any $0<\varepsilon\leq1/2$, the quantity $\varepsilon/8$ is positive.
\end{theorem}

\begin{definition}[Learning consequences degree]
\lean{FABL.learningConsequencesDegree}
\statementsource{lmn-constant-depth-fourier}{pdf-63759bd909d5-p004-b006}
\statementsource{Analysis-of-Boolean-Functions-by-Ryan-ODonnell}{
  pdf-ffbe9c7b3306-p069-b004,
  pdf-ffbe9c7b3306-p081-b005,
  pdf-ffbe9c7b3306-p268-b003
}
\label{FABL.learningConsequencesDegree}
\leanok
\uses{FABL.PositiveInfluenceBound, FABL.PositiveLearningParameter, FABL.learningConsequencesInternalError}
Let $t>0$ be a rational influence bound and let $0<\varepsilon\leq 1/2$ be a rational
learning parameter. The learning consequences degree is the ceiling of $t$ divided by
the internal error associated with $\varepsilon$.
\end{definition}

\begin{definition}[Learning consequences threshold]
\lean{FABL.learningConsequencesThreshold}
\statementsource{lmn-constant-depth-fourier}{pdf-63759bd909d5-p009-b001}
\statementsource{Analysis-of-Boolean-Functions-by-Ryan-ODonnell}{
  pdf-ffbe9c7b3306-p069-b003,
  pdf-ffbe9c7b3306-p069-b004,
  pdf-ffbe9c7b3306-p268-b003
}
\label{FABL.learningConsequencesThreshold}
\leanok
\uses{FABL.PositiveInfluenceBound, FABL.PositiveLearningParameter, FABL.learningConsequencesDegree, FABL.learningConsequencesInternalError}
Let $t>0$ be a rational influence bound and let $0<\varepsilon\leq 1/2$ be a rational
learning parameter. The learning consequences threshold is the internal error associated
with $\varepsilon$, squared and divided by $t^2$, multiplied by $(1/9)^d$, where $d$ is
the learning consequences degree for $t$ and $\varepsilon$.
\end{definition}

\begin{theorem}[Positivity of the threshold]
\lean{FABL.learningConsequencesThreshold_pos}
\label{FABL.learningConsequencesThreshold_pos}
\leanok
\uses{FABL.PositiveInfluenceBound, FABL.PositiveLearningParameter, FABL.learningConsequencesInternalError_pos, FABL.learningConsequencesThreshold}
\difficulty{3}
For any $t>0$ and any $0<\varepsilon\leq1/2$, define $d=\lceil8t/\varepsilon\rceil$ and
$\theta=(\varepsilon/8)^2t^{-2}(1/9)^d$. Then $\theta>0$.
\end{theorem}

\begin{definition}[Learning consequences coordinate cap]
\lean{FABL.learningConsequencesCoordinateCap}
\statementsource{lmn-constant-depth-fourier}{pdf-63759bd909d5-p010-b002}
\statementsource{Analysis-of-Boolean-Functions-by-Ryan-ODonnell}{
  pdf-ffbe9c7b3306-p071-b002,
  pdf-ffbe9c7b3306-p071-b004,
  pdf-ffbe9c7b3306-p268-b003
}
\label{FABL.learningConsequencesCoordinateCap}
\leanok
\uses{FABL.PositiveInfluenceBound, FABL.PositiveLearningParameter, FABL.learningConsequencesThreshold}
Let $t>0$ be a rational influence bound and let $0<\varepsilon\leq 1/2$ be a rational
learning parameter. The learning consequences coordinate cap is the ceiling of $2t$
divided by the learning consequences threshold for $t$ and $\varepsilon$.
\end{definition}

\begin{definition}[Learning consequences family bound]
\lean{FABL.learningConsequencesFamilyBound}
\statementsource{lmn-constant-depth-fourier}{pdf-63759bd909d5-p009-b001}
\statementsource{Analysis-of-Boolean-Functions-by-Ryan-ODonnell}{
  pdf-ffbe9c7b3306-p079-b003,
  pdf-ffbe9c7b3306-p081-b003,
  pdf-ffbe9c7b3306-p268-b003
}
\label{FABL.learningConsequencesFamilyBound}
\leanok
\uses{FABL.PositiveInfluenceBound, FABL.PositiveLearningParameter, FABL.learningConsequencesCoordinateCap, FABL.learningConsequencesDegree}
Let $t>0$ be a rational influence bound and let $0<\varepsilon\leq 1/2$ be a rational
learning parameter. If $d$ is the learning consequences degree and $c$ is the learning
consequences coordinate cap for these parameters, the family bound is $(d+1)(c+1)^d$.
\end{definition}

\begin{theorem}[Positivity of the family bound]
\lean{FABL.learningConsequencesFamilyBound_pos}
\label{FABL.learningConsequencesFamilyBound_pos}
\leanok
\uses{FABL.PositiveInfluenceBound, FABL.PositiveLearningParameter, FABL.learningConsequencesFamilyBound}
\difficulty{3}
For any $t>0$ and any $0<\varepsilon\leq1/2$, define $d=\lceil8t/\varepsilon\rceil$,
$\theta=(\varepsilon/8)^2t^{-2}(1/9)^d$, and $k=\lceil2t/\theta\rceil$. Then
$(d+1)(k+1)^d>0$.
\end{theorem}

\begin{definition}[Learning consequences target family]
\lean{FABL.learningConsequencesTargetFamily}
\statementsource{lmn-constant-depth-fourier}{
  pdf-63759bd909d5-p010-b001,
  pdf-63759bd909d5-p010-b004
}
\statementsource{Analysis-of-Boolean-Functions-by-Ryan-ODonnell}{
  pdf-ffbe9c7b3306-p079-b003,
  pdf-ffbe9c7b3306-p268-b003
}
\label{FABL.learningConsequencesTargetFamily}
\leanok
\uses{FABL.BooleanFunction, FABL.PositiveInfluenceBound, FABL.PositiveLearningParameter, FABL.influentialCoordinates, FABL.learningConsequencesDegree, FABL.learningConsequencesThreshold, FABL.lowDegreeSubsets}
Let $n$ be a natural number, let $f:\{−1,1\}^n\to\{−1,1\}$ be a Boolean function, let
$t>0$ be a rational influence bound, and let $0<\varepsilon\leq 1/2$ be a rational
learning parameter. The target family consists of the subsets of the influential
coordinates of $f$ whose cardinalities are at most the learning consequences degree for
$t$ and $\varepsilon$, with influence measured at the learning consequences threshold
for those parameters.
\end{definition}

\begin{theorem}[Exterior weight for influential low-degree subsets]
\lean{FABL.fourierWeightOutside_lowDegreeSubsets_influential_le_add}
\label{FABL.fourierWeightOutside_lowDegreeSubsets_influential_le_add}
\leanok
\uses{FABL.BooleanFunction, FABL.BooleanFunction.toReal, FABL.fourierWeight, FABL.fourierWeightAbove, FABL.fourierWeightOutside, FABL.influentialCoordinates, FABL.lowDegreeOutsideInfluentialMass, FABL.lowDegreeSubsets}
\difficulty{5}
Let $n$ and $k$ be natural numbers, let $f$ be a sign-valued function on the
$n$-dimensional sign cube, and let $\tau\in\mathbb{R}$. The sum of the squared Fourier
coefficients of its real-valued version over sets outside the family of subsets of size
at most $k$ contained in the coordinates with Boolean influence at least $\tau$ is at
most the sum of squared Fourier coefficients over sets of size at most $k$ that are not
contained in those coordinates, plus the sum over sets of size greater than $k$.
\end{theorem}

\begin{theorem}[Square Root of the Learning Threshold]
\lean{FABL.sqrt_learningConsequencesThreshold}
\label{FABL.sqrt_learningConsequencesThreshold}
\leanok
\uses{FABL.PositiveInfluenceBound, FABL.PositiveLearningParameter, FABL.learningConsequencesDegree, FABL.learningConsequencesInternalError, FABL.learningConsequencesInternalError_pos, FABL.learningConsequencesThreshold}
\difficulty{5}
Let $t$ be a positive rational number and let $\varepsilon$ be rational with
$0<\varepsilon\le \tfrac12$. Set $e=\varepsilon/8$, $d=\lceil t/e\rceil$, and
$\tau=(e^2/t^2)(1/9)^d$. Then $\sqrt{\tau}= (e/t)(1/3)^d$, with the quantities on both
sides regarded as real numbers.
\end{theorem}

\begin{theorem}[Fourier concentration at a learning degree]
\lean{FABL.spectrum_concentrated_up_to_learningConsequencesDegree}
\proofsource{Analysis-of-Boolean-Functions-by-Ryan-ODonnell}{
  pdf-ffbe9c7b3306-p069-b003,
  pdf-ffbe9c7b3306-p069-b004,
  pdf-ffbe9c7b3306-p268-b003,
  pdf-ffbe9c7b3306-p268-b004,
  pdf-ffbe9c7b3306-p269-b001
}
\label{FABL.spectrum_concentrated_up_to_learningConsequencesDegree}
\leanok
\uses{FABL.BooleanFunction, FABL.BooleanFunction.toReal, FABL.IsFourierSpectrumConcentratedUpTo, FABL.IsFourierSpectrumConcentratedUpTo.mono_cutoff, FABL.PositiveInfluenceBound, FABL.PositiveLearningParameter, FABL.isFourierSpectrumConcentratedUpTo_totalInfluence_div, FABL.learningConsequencesDegree, FABL.learningConsequencesInternalError, FABL.learningConsequencesInternalError_pos, FABL.totalInfluence}
\difficulty{4}
Let $f:\{\pm1\}^n\to\{\pm1\}$, let $t$ be a positive rational number, and let
$\varepsilon$ be rational with $0<\varepsilon\leq 1/2$. If the total influence of $f$ is
at most $t$, then the sum of its squared Fourier coefficients over subsets whose size is
greater than $\lceil 8t/\varepsilon\rceil$ is at most $\varepsilon/8$.
\end{theorem}

\begin{theorem}[Learning consequences target family spectrum concentrated]
\lean{FABL.learningConsequencesTargetFamily_spectrum_concentrated}
\statementsource{lmn-constant-depth-fourier}{
  pdf-63759bd909d5-p002-b004,
  pdf-63759bd909d5-p009-b001
}
\proofsource{Analysis-of-Boolean-Functions-by-Ryan-ODonnell}{
  pdf-ffbe9c7b3306-p268-b003,
  pdf-ffbe9c7b3306-p268-b004,
  pdf-ffbe9c7b3306-p269-b001
}
\label{FABL.learningConsequencesTargetFamily_spectrum_concentrated}
\leanok
\uses{FABL.BooleanFunction, FABL.BooleanFunction.toReal, FABL.IsFourierSpectrumConcentratedOn, FABL.IsFourierSpectrumConcentratedUpTo, FABL.PositiveInfluenceBound, FABL.PositiveLearningParameter, FABL.fourierWeightAbove, FABL.fourierWeightAboveReal_natCast, FABL.fourierWeightOutside, FABL.fourierWeightOutside_lowDegreeSubsets_influential_le_add, FABL.influentialCoordinates, FABL.learningConsequencesDegree, FABL.learningConsequencesInternalError, FABL.learningConsequencesInternalError_pos, FABL.learningConsequencesTargetFamily, FABL.learningConsequencesThreshold, FABL.learningConsequencesThreshold_pos, FABL.lowDegreeOutsideInfluentialMass, FABL.lowDegreeOutsideInfluentialMass_le, FABL.lowDegreeSubsets, FABL.spectrum_concentrated_up_to_learningConsequencesDegree, FABL.sqrt_learningConsequencesThreshold, FABL.totalInfluence}
\difficulty{5}
% dataset-pipeline: stub -- prose pending build_informal_statements.py
Statement pending regeneration.
\end{theorem}

\begin{theorem}[Card influential coordinates learning consequences threshold le]
\lean{FABL.card_influentialCoordinates_learningConsequencesThreshold_le}
\statementsource{Analysis-of-Boolean-Functions-by-Ryan-ODonnell}{pdf-ffbe9c7b3306-p071-b004}
\label{FABL.card_influentialCoordinates_learningConsequencesThreshold_le}
\leanok
\uses{FABL.BooleanFunction, FABL.BooleanFunction.toReal, FABL.PositiveInfluenceBound, FABL.PositiveLearningParameter, FABL.card_influentialCoordinates_le_div, FABL.influentialCoordinates, FABL.learningConsequencesCoordinateCap, FABL.learningConsequencesThreshold, FABL.learningConsequencesThreshold_pos, FABL.totalInfluence}
\difficulty{4}
% dataset-pipeline: stub -- prose pending build_informal_statements.py
Statement pending regeneration.
\end{theorem}

\begin{theorem}[Learning target family size bound]
\lean{FABL.card_learningConsequencesTargetFamily_le}
\label{FABL.card_learningConsequencesTargetFamily_le}
\leanok
\uses{FABL.BooleanFunction, FABL.BooleanFunction.toReal, FABL.PositiveInfluenceBound, FABL.PositiveLearningParameter, FABL.card_influentialCoordinates_learningConsequencesThreshold_le, FABL.card_lowDegreeSubsets_le, FABL.influentialCoordinates, FABL.learningConsequencesCoordinateCap, FABL.learningConsequencesDegree, FABL.learningConsequencesFamilyBound, FABL.learningConsequencesTargetFamily, FABL.learningConsequencesThreshold, FABL.totalInfluence}
\difficulty{4}
Let $n$ be a natural number, let $f:\{-1,1\}^n\to\{\pm1\}$, let $t\in\mathbb Q$ be
positive, and let $\varepsilon\in\mathbb Q$ satisfy $0<\varepsilon\leq 1/2$. If the
total influence of the real-valued version of $f$ is at most $t$, then the
learning-consequences target family determined by $f,t,\varepsilon$ has cardinality at
most the learning-consequences family bound determined by $t,\varepsilon$.
\end{theorem}

\begin{definition}[Bounded influence query learning program]
\lean{FABL.boundedInfluenceQueryLearningProgram}
\statementsource{lmn-constant-depth-fourier}{
  pdf-63759bd909d5-p010-b001,
  pdf-63759bd909d5-p010-b002,
  pdf-63759bd909d5-p010-b003
}
\statementsource{Analysis-of-Boolean-Functions-by-Ryan-ODonnell}{
  pdf-ffbe9c7b3306-p079-b001,
  pdf-ffbe9c7b3306-p081-b005,
  pdf-ffbe9c7b3306-p273-b008
}
\label{FABL.boundedInfluenceQueryLearningProgram}
\leanok
\uses{FABL.LearningProgram, FABL.PositiveInfluenceBound, FABL.PositiveLearningParameter, FABL.SparseFourierHypothesis, FABL.kushilevitzMansourProgramForBound, FABL.learningConsequencesFamilyBound}
Let $n$ be a natural number, let $t>0$ be a rational influence bound, and let
$0<\varepsilon\leq 1/2$ be a rational learning parameter. The bounded influence query
learning program is a membership-query learning program on the $n$-dimensional sign cube
that outputs a sparse Fourier hypothesis, obtained by applying the Kushilevitz–Mansour
program with accuracy parameter $\varepsilon$ and the learning consequences family bound
for $t$ and $\varepsilon$.
\end{definition}

\begin{theorem}[Bounded influence query learning program failure probability le one tenth]
\lean{FABL.boundedInfluenceQueryLearningProgram_failureProbability_le_one_tenth}
\statementsource{Analysis-of-Boolean-Functions-by-Ryan-ODonnell}{
  pdf-ffbe9c7b3306-p079-b001,
  pdf-ffbe9c7b3306-p081-b005,
  pdf-ffbe9c7b3306-p273-b008
}
\label{FABL.boundedInfluenceQueryLearningProgram_failureProbability_le_one_tenth}
\leanok
\uses{FABL.BooleanFunction, FABL.BooleanFunction.toReal, FABL.LearningProgram.eventProbability, FABL.PositiveInfluenceBound, FABL.PositiveLearningParameter, FABL.SparseFourierHypothesis.evaluate, FABL.boundedInfluenceQueryLearningProgram, FABL.card_learningConsequencesTargetFamily_le, FABL.kushilevitzMansourProgramForBound_failureProbability_le_one_tenth, FABL.learningConsequencesFamilyBound, FABL.learningConsequencesTargetFamily, FABL.learningConsequencesTargetFamily_spectrum_concentrated, FABL.relativeHammingDist, FABL.totalInfluence}
\difficulty{1}
% dataset-pipeline: stub -- prose pending build_informal_statements.py
Statement pending regeneration.
\end{theorem}

\begin{theorem}[Query bound for bounded influence learning]
\lean{FABL.boundedInfluenceQueryLearningProgram_queries_le}
\label{FABL.boundedInfluenceQueryLearningProgram_queries_le}
\leanok
\uses{FABL.BooleanFunction, FABL.LearningCost, FABL.LearningProgram.runWithCost, FABL.PositiveInfluenceBound, FABL.PositiveLearningParameter, FABL.SparseFourierHypothesis, FABL.boundedInfluenceQueryLearningProgram, FABL.kushilevitzMansourProgramForBound_queries_polynomial_cast_le, FABL.learningConsequencesFamilyBound}
\difficulty{1}
Let $n$ be a nonnegative integer, let $f$ be a Boolean function on the $n$-dimensional
sign cube, let $t$ be a positive rational influence bound, and let $\epsilon$ be a
rational learning parameter with $0<\epsilon\le 1/2$. For every hypothesis and cost pair
in the support of the query-learning program run with target $f$, the number of
membership queries is at most $2^{40}(n+1)^2(B+1)^8/\epsilon^{10}$, where $B$ is the
family bound determined by $t$ and $\epsilon$.
\end{theorem}

\begin{theorem}[Bounded influence query learning program work le]
\lean{FABL.boundedInfluenceQueryLearningProgram_work_le}
\statementsource{lmn-constant-depth-fourier}{pdf-63759bd909d5-p010-b001}
\label{FABL.boundedInfluenceQueryLearningProgram_work_le}
\leanok
\uses{FABL.BooleanFunction, FABL.LearningCost, FABL.LearningProgram.runWithCost, FABL.PositiveInfluenceBound, FABL.PositiveLearningParameter, FABL.SparseFourierHypothesis, FABL.boundedInfluenceQueryLearningProgram, FABL.kushilevitzMansourProgramForBound_work_polynomial_cast_le, FABL.learningConsequencesFamilyBound}
\difficulty{1}
% dataset-pipeline: stub -- prose pending build_informal_statements.py
Statement pending regeneration.
\end{theorem}

\begin{definition}[Singleton Fourier family]
\lean{FABL.singletonFourierFamily}
\label{FABL.singletonFourierFamily}
\leanok
For each natural number $n$, the singleton Fourier family consists of all one-element
subsets of the coordinate set $\{0,\ldots,n-1\}$.
\end{definition}

\begin{theorem}[Singletons in the singleton Fourier family]
\lean{FABL.singleton_mem_singletonFourierFamily}
\label{FABL.singleton_mem_singletonFourierFamily}
\leanok
\uses{FABL.singletonFourierFamily}
\difficulty{1}
For every nonnegative integer $n$ and every coordinate $i$ in $\{0,\dots,n-1\}$, the
singleton set $\{i\}$ belongs to the family of all singleton subsets of the coordinate
set $\{0,\dots,n-1\}$.
\end{theorem}

\begin{theorem}[Size of the singleton Fourier family]
\lean{FABL.card_singletonFourierFamily}
\label{FABL.card_singletonFourierFamily}
\leanok
\uses{FABL.singletonFourierFamily}
\difficulty{3}
For every natural number $n$, the family of singleton subsets of the $n$ coordinate
indices has cardinality $n$.
\end{theorem}

\begin{definition}[Singleton influence accuracy]
\lean{FABL.singletonInfluenceAccuracy}
\statementsource{lmn-constant-depth-fourier}{pdf-63759bd909d5-p010-b005}
\statementsource{Analysis-of-Boolean-Functions-by-Ryan-ODonnell}{pdf-ffbe9c7b3306-p080-b002}
\label{FABL.singletonInfluenceAccuracy}
\leanok
\uses{FABL.PositiveInfluenceBound, FABL.PositiveLearningParameter, FABL.learningConsequencesThreshold, FABL.learningConsequencesThreshold_pos}
Let $t>0$ be a rational influence bound and let $0<\varepsilon\leq 1/2$ be a rational
learning parameter. The singleton influence accuracy is the smaller of one half and one
quarter of the learning consequences threshold for $t$ and $\varepsilon$.
\end{definition}

\begin{theorem}[Singleton accuracy value]
\lean{FABL.singletonInfluenceAccuracy_value}
\label{FABL.singletonInfluenceAccuracy_value}
\leanok
\uses{FABL.PositiveInfluenceBound, FABL.PositiveLearningParameter, FABL.learningConsequencesThreshold, FABL.singletonInfluenceAccuracy}
\difficulty{1}
Let $t$ be a positive rational and let $\varepsilon$ be rational with $0<\varepsilon\le
1/2$. Define $D=\lceil 8t/\varepsilon\rceil$ and $T=(\varepsilon/8)^2t^{-2}(1/9)^D$. The
singleton influence accuracy is exactly $\min(T/4,1/2)$.
\end{theorem}

\begin{theorem}[Singleton accuracy threshold bound]
\lean{FABL.singletonInfluenceAccuracy_le_quarter_threshold}
\label{FABL.singletonInfluenceAccuracy_le_quarter_threshold}
\leanok
\uses{FABL.PositiveInfluenceBound, FABL.PositiveLearningParameter, FABL.learningConsequencesThreshold, FABL.singletonInfluenceAccuracy}
\difficulty{1}
Let $t$ be a positive rational and let $\varepsilon$ be rational with $0<\varepsilon\le
1/2$. Define $D=\lceil 8t/\varepsilon\rceil$ and $T=(\varepsilon/8)^2t^{-2}(1/9)^D$.
Then the singleton influence accuracy $\min(T/4,1/2)$ is at most $T/4$.
\end{theorem}

\begin{definition}[Singleton influence confidence]
\lean{FABL.singletonInfluenceConfidence}
\statementsource{lmn-constant-depth-fourier}{pdf-63759bd909d5-p010-b005}
\statementsource{Analysis-of-Boolean-Functions-by-Ryan-ODonnell}{pdf-ffbe9c7b3306-p080-b002}
\label{FABL.singletonInfluenceConfidence}
\leanok
\uses{FABL.PositiveLearningParameter}
For each natural number $n$, the singleton influence confidence is the rational number
$1/(20(n+1))$.
\end{definition}

\begin{definition}[Singleton Influence Sample Count]
\lean{FABL.singletonInfluenceSampleCount}
\statementsource{lmn-constant-depth-fourier}{pdf-63759bd909d5-p010-b002}
\statementsource{Analysis-of-Boolean-Functions-by-Ryan-ODonnell}{
  pdf-ffbe9c7b3306-p080-b002,
  pdf-ffbe9c7b3306-p080-b003
}
\label{FABL.singletonInfluenceSampleCount}
\leanok
\uses{FABL.PositiveInfluenceBound, FABL.PositiveLearningParameter, FABL.fourierEstimatorSampleCount, FABL.singletonInfluenceAccuracy, FABL.singletonInfluenceConfidence}
For a dimension $n\in\mathbb{N}$, a positive rational influence bound $t$, and a
rational learning parameter $0<\varepsilon\leq 1/2$, define the singleton influence
sample count to be the Fourier estimator sample count at accuracy prescribed by $t$ and
$\varepsilon$, and confidence prescribed by $n$.
\end{definition}

\begin{definition}[Uncapped empirical influential coordinates]
\lean{FABL.empiricalInfluentialCoordinatesUncapped}
\label{FABL.empiricalInfluentialCoordinatesUncapped}
\leanok
\uses{FABL.PositiveInfluenceBound, FABL.PositiveLearningParameter, FABL.Sign, FABL.empiricalFourierCoeff, FABL.learningConsequencesThreshold, FABL.singletonFourierFamily, FABL.singletonInfluenceAccuracy, FABL.singletonInfluenceSampleCount, FABL.singleton_mem_singletonFourierFamily}
Let $n$ be a natural number, let $t$ be a positive rational influence bound, and let
$\varepsilon$ be a rational with $0<\varepsilon\le 1/2$. Given, for each singleton
coordinate set, a list of the prescribed sample length of pairs consisting of a sign
vector of length $n$ and a sign, include a coordinate $i$ in the empirical set exactly
when its empirical Fourier coefficient, computed from the samples assigned to $\{i\}$,
is at least the learning threshold for $t,\varepsilon$ minus the prescribed accuracy for
$t,\varepsilon$.
\end{definition}

\begin{definition}[Capped empirical influential coordinates]
\lean{FABL.empiricalInfluentialCoordinates}
\label{FABL.empiricalInfluentialCoordinates}
\leanok
\uses{FABL.PositiveInfluenceBound, FABL.PositiveLearningParameter, FABL.Sign, FABL.empiricalInfluentialCoordinatesUncapped, FABL.learningConsequencesCoordinateCap, FABL.singletonFourierFamily, FABL.singletonInfluenceSampleCount}
Let $n$ be a natural number, let $t$ be a positive rational influence bound, and let
$\varepsilon$ be a rational with $0<\varepsilon\le 1/2$. Form the uncapped empirical set
using the prescribed samples; the empirical set is that set if its cardinality is at
most the coordinate cap for $t,\varepsilon$, and is empty otherwise.
\end{definition}

\begin{theorem}[Bound on selected coordinates]
\lean{FABL.card_empiricalInfluentialCoordinates_le}
\label{FABL.card_empiricalInfluentialCoordinates_le}
\leanok
\uses{FABL.PositiveInfluenceBound, FABL.PositiveLearningParameter, FABL.Sign, FABL.empiricalInfluentialCoordinates, FABL.empiricalInfluentialCoordinatesUncapped, FABL.learningConsequencesCoordinateCap, FABL.singletonFourierFamily, FABL.singletonInfluenceSampleCount}
\difficulty{4}
For a positive rational influence bound $t$, a learning parameter $\varepsilon$ with
$0<\varepsilon\leq 1/2$, and any prescribed samples for every coordinate, the empirical
influential-coordinate set has cardinality at most its coordinate cap, defined from $t$
and $\varepsilon$.
\end{theorem}

\begin{definition}[Empirical influential coordinates from inputs]
\lean{FABL.empiricalInfluentialCoordinatesFromInputs}
\label{FABL.empiricalInfluentialCoordinatesFromInputs}
\leanok
\uses{FABL.BooleanFunction, FABL.PositiveInfluenceBound, FABL.PositiveLearningParameter, FABL.empiricalInfluentialCoordinates, FABL.singletonFourierFamily, FABL.singletonInfluenceSampleCount}
Let $n$ be a natural number, let $f$ be a Boolean function from length-$n$ sign vectors
to signs, let $t$ be a positive rational influence bound, and let $\varepsilon$ be a
rational with $0<\varepsilon\le 1/2$. Given, for every singleton coordinate set, a
prescribed-length list of input sign vectors, form the empirical influential-coordinate
set using the sample pairs consisting of each input and its value under $f$.
\end{definition}

\begin{definition}[Monotone Influence Selection Program]
\lean{FABL.monotoneInfluenceSelectionProgram}
\statementsource{Analysis-of-Boolean-Functions-by-Ryan-ODonnell}{pdf-ffbe9c7b3306-p273-b008}
\label{FABL.monotoneInfluenceSelectionProgram}
\leanok
\uses{FABL.LearningProgram, FABL.PositiveInfluenceBound, FABL.PositiveLearningParameter, FABL.empiricalInfluentialCoordinates, FABL.singletonFourierFamily, FABL.singletonInfluenceSampleCount}
For a dimension $n\in\mathbb{N}$, a positive rational influence bound $t$, and a
rational learning parameter $0<\varepsilon\leq 1/2$, define a random-example learning
program that draws a matrix of independent uniform labeled examples, with one sample
batch for each member of the singleton Fourier family. The batch size is the singleton
influence sample count, the charged local work is the family size times the batch size
times $n+1$, plus $n$, and the output is the set of empirically influential coordinates.
\end{definition}

\begin{theorem}[Exact cost of influence selection]
\lean{FABL.monotoneInfluenceSelectionProgram_cost_eq}
\label{FABL.monotoneInfluenceSelectionProgram_cost_eq}
\leanok
\uses{FABL.BooleanFunction, FABL.LearningCost, FABL.LearningProgram.randomExampleMatrixOutcome, FABL.LearningProgram.randomExampleMatrixOutputLaw, FABL.LearningProgram.runWithCost, FABL.PositiveInfluenceBound, FABL.PositiveLearningParameter, FABL.card_singletonFourierFamily, FABL.monotoneInfluenceSelectionProgram, FABL.singletonInfluenceSampleCount}
\difficulty{4}
Let $n$ be a nonnegative integer, let $t$ be a positive rational number, and let
$0<\varepsilon\le 1/2$. For any Boolean target function and any outcome in the support
of the influence-selection program with parameters $n,t,\varepsilon$, the cost is
exactly $\bigl(nm,\,0,\,nm+nm(n+1)+n\bigr)$, where $m$ is the program's
singleton-influence sample count for these parameters.
\end{theorem}

\begin{theorem}[Bound on low degree subsets of selected coordinates]
\lean{FABL.card_lowDegreeSubsets_empiricalInfluentialCoordinates_le}
\label{FABL.card_lowDegreeSubsets_empiricalInfluentialCoordinates_le}
\leanok
\uses{FABL.PositiveInfluenceBound, FABL.PositiveLearningParameter, FABL.Sign, FABL.card_empiricalInfluentialCoordinates_le, FABL.card_lowDegreeSubsets_le, FABL.empiricalInfluentialCoordinates, FABL.learningConsequencesCoordinateCap, FABL.learningConsequencesDegree, FABL.learningConsequencesFamilyBound, FABL.lowDegreeSubsets, FABL.singletonFourierFamily, FABL.singletonInfluenceSampleCount}
\difficulty{3}
Let $n$ be a nonnegative integer, let $t$ be a positive rational number, and let
$\varepsilon$ be a rational number with $0<\varepsilon\le 1/2$. From any assignment of
samples to each singleton coordinate set, form the empirically influential coordinates
using the given threshold and cap. The number of subsets of those coordinates of size at
most the degree determined by $t$ and $\varepsilon$ is at most the corresponding family
bound.
\end{theorem}

\begin{theorem}[Bound on half-threshold influential coordinates]
\lean{FABL.card_halfThresholdInfluentialCoordinates_le}
\label{FABL.card_halfThresholdInfluentialCoordinates_le}
\leanok
\uses{FABL.BooleanFunction, FABL.BooleanFunction.toReal, FABL.PositiveInfluenceBound, FABL.PositiveLearningParameter, FABL.card_influentialCoordinates_le_div, FABL.influentialCoordinates, FABL.learningConsequencesCoordinateCap, FABL.learningConsequencesThreshold, FABL.learningConsequencesThreshold_pos, FABL.totalInfluence}
\difficulty{4}
Let $n$ be a natural number and let $f:\{-1,1\}^n\to\{-1,1\}$ be a Boolean
function. Let $t$ be a positive rational and let $\varepsilon$ be a rational with
$0<\varepsilon\leq\frac12$. For each coordinate $i$, let its influence be the uniform
probability that flipping $i$ changes $f$, and let the total influence be the sum of the
squared discrete derivatives of the associated real-valued function. If this total
influence is at most $t$, then the number of coordinates whose influence is at least
half of $\theta=\frac12(\varepsilon/8)^2t^{-2}(1/9)^{\lceil t/
(\varepsilon/8)\rceil}$ is at most $\left\lceil 2t/\theta\right\rceil$.
\end{theorem}

\begin{theorem}[Singleton Fourier coefficient and influence]
\lean{FABL.fourierCoeff_singleton_eq_booleanInfluence_of_monotone}
\proofsource{Analysis-of-Boolean-Functions-by-Ryan-ODonnell}{
  pdf-ffbe9c7b3306-p048-b002,
  pdf-ffbe9c7b3306-p048-b003
}
\label{FABL.fourierCoeff_singleton_eq_booleanInfluence_of_monotone}
\leanok
\uses{FABL.BooleanFunction, FABL.BooleanFunction.toReal, FABL.booleanInfluence, FABL.booleanInfluence_eq_influence_toReal, FABL.fourierCoeff, FABL.influence_eq_fourierCoeff_singleton_of_monotone}
\difficulty{1}
Let $n$ be a natural number, let $f:\{-1,1\}^n\to\{-1,1\}$ be monotone in the
coordinatewise order, and choose a coordinate $i\in\{1,\ldots,n\}$. Then the Fourier
coefficient of $f$ at the singleton $\{i\}$ equals the probability, under the uniform
distribution on the cube, that flipping coordinate $i$ changes the value of $f$.
\end{theorem}

\begin{theorem}[Empirical detection of influential coordinates]
\lean{FABL.influentialCoordinates_subset_empiricalInfluentialCoordinates_of_no_bad}
\label{FABL.influentialCoordinates_subset_empiricalInfluentialCoordinates_of_no_bad}
\leanok
\uses{FABL.BooleanFunction, FABL.BooleanFunction.toReal, FABL.PositiveInfluenceBound, FABL.PositiveLearningParameter, FABL.Sign, FABL.booleanInfluence, FABL.card_halfThresholdInfluentialCoordinates_le, FABL.empiricalFourierCoeff, FABL.empiricalFourierCoeff_cast, FABL.empiricalInfluentialCoordinates, FABL.empiricalInfluentialCoordinatesFromInputs, FABL.empiricalInfluentialCoordinatesUncapped, FABL.finiteFamilyCoefficientBadSetWithParameters, FABL.fourierCoeff, FABL.fourierCoeff_singleton_eq_booleanInfluence_of_monotone, FABL.influentialCoordinates, FABL.learningConsequencesCoordinateCap, FABL.learningConsequencesThreshold, FABL.mem_influentialCoordinates, FABL.realEmpiricalFourierCoeff, FABL.singletonFourierFamily, FABL.singletonInfluenceAccuracy, FABL.singletonInfluenceAccuracy_le_quarter_threshold, FABL.singletonInfluenceConfidence, FABL.singletonInfluenceSampleCount, FABL.singleton_mem_singletonFourierFamily, FABL.totalInfluence}
\difficulty{5}
Let $n$ be a nonnegative integer and let $f:\{−1,1\}^n\to\{−1,1\}$ be monotone. Let
$t\in\mathbb Q$ be positive and let $\varepsilon\in\mathbb Q$ satisfy $0<\varepsilon\le
1/2$; suppose the total influence of the real-valued sign version of $f$ is at most $t$.
For each singleton Fourier query, take the prescribed sample of inputs, and suppose its
empirical Fourier coefficient differs from the true coefficient by less than the
prescribed accuracy. Then every coordinate whose Boolean influence is at least the
prescribed learning threshold belongs to the empirical set of influential coordinates
formed from these samples.
\end{theorem}

\begin{definition}[Bad monotone-selection condition]
\lean{FABL.monotoneSelectionBad}
\label{FABL.monotoneSelectionBad}
\leanok
\uses{FABL.BooleanFunction, FABL.PositiveInfluenceBound, FABL.PositiveLearningParameter, FABL.influentialCoordinates, FABL.learningConsequencesCoordinateCap, FABL.learningConsequencesThreshold}
Let $n$ be a natural number, let $f$ be a Boolean function from length-$n$ sign vectors
to signs, let $t$ be a positive rational influence bound, let $\varepsilon$ be a
rational with $0<\varepsilon\le 1/2$, and let $J$ be a set of coordinates. The selection
is bad if either the set of coordinates whose influence under $f$ is at least the
learning threshold for $t,\varepsilon$ is not contained in $J$, or the cardinality of
$J$ exceeds the coordinate cap for $t,\varepsilon$.
\end{definition}

\begin{theorem}[Monotone influence selection program failure probability le one twentieth]
\lean{FABL.monotoneInfluenceSelectionProgram_failureProbability_le_one_twentieth}
\statementsource{Analysis-of-Boolean-Functions-by-Ryan-ODonnell}{pdf-ffbe9c7b3306-p273-b008}
\label{FABL.monotoneInfluenceSelectionProgram_failureProbability_le_one_twentieth}
\leanok
\uses{FABL.BooleanFunction, FABL.BooleanFunction.toReal, FABL.LearningProgram.eventProbability, FABL.LearningProgram.randomExampleMatrixInputLaw, FABL.LearningProgram.randomExampleMatrixOutcome, FABL.LearningProgram.randomExampleMatrixOutputLaw, FABL.LearningProgram.randomExampleMatrixOutputLaw_toOuterMeasure_apply, FABL.LearningProgram.runWithCost, FABL.PositiveInfluenceBound, FABL.PositiveLearningParameter, FABL.card_empiricalInfluentialCoordinates_le, FABL.card_singletonFourierFamily, FABL.empiricalInfluentialCoordinates, FABL.empiricalInfluentialCoordinatesFromInputs, FABL.finiteFamilyCoefficientBadSetWithParameters, FABL.influentialCoordinates_subset_empiricalInfluentialCoordinates_of_no_bad, FABL.learningConsequencesCoordinateCap, FABL.measure_finiteFamily_someCoefficientBadSetWithParameters_le, FABL.monotoneInfluenceSelectionProgram, FABL.monotoneSelectionBad, FABL.singletonFourierFamily, FABL.singletonInfluenceAccuracy, FABL.singletonInfluenceConfidence, FABL.singletonInfluenceSampleCount, FABL.totalInfluence}
\difficulty{6}
% dataset-pipeline: stub -- prose pending build_informal_statements.py
Statement pending regeneration.
\end{theorem}

\begin{definition}[Bounded Family Coefficient Accuracy]
\lean{FABL.boundedFamilyCoefficientAccuracy}
\statementsource{lmn-constant-depth-fourier}{pdf-63759bd909d5-p010-b005}
\statementsource{Analysis-of-Boolean-Functions-by-Ryan-ODonnell}{pdf-ffbe9c7b3306-p080-b002}
\label{FABL.boundedFamilyCoefficientAccuracy}
\leanok
\uses{FABL.PositiveLearningParameter}
Let $0<\varepsilon\leq 1/2$ be a rational learning parameter, and let $M$ be a positive
integer. The bounded-family coefficient accuracy is the rational parameter
$\varepsilon/(2M)$.
\end{definition}

\begin{definition}[Bounded Family Coefficient Confidence]
\lean{FABL.boundedFamilyCoefficientConfidence}
\statementsource{lmn-constant-depth-fourier}{pdf-63759bd909d5-p010-b005}
\statementsource{Analysis-of-Boolean-Functions-by-Ryan-ODonnell}{pdf-ffbe9c7b3306-p080-b002}
\label{FABL.boundedFamilyCoefficientConfidence}
\leanok
\uses{FABL.PositiveLearningParameter}
Let $M$ be a positive integer. The bounded-family coefficient confidence is the rational
parameter $1/(20M)$.
\end{definition}

\begin{definition}[Bounded Family Sample Count]
\lean{FABL.boundedFamilySampleCount}
\statementsource{lmn-constant-depth-fourier}{pdf-63759bd909d5-p010-b002}
\statementsource{Analysis-of-Boolean-Functions-by-Ryan-ODonnell}{
  pdf-ffbe9c7b3306-p080-b002,
  pdf-ffbe9c7b3306-p080-b003
}
\label{FABL.boundedFamilySampleCount}
\leanok
\uses{FABL.PositiveLearningParameter, FABL.boundedFamilyCoefficientAccuracy, FABL.boundedFamilyCoefficientConfidence, FABL.fourierEstimatorSampleCount}
Let $0<\varepsilon\leq 1/2$ be a rational learning parameter, and let $M$ be a positive
integer. Define the bounded-family sample count as the Fourier estimator sample count at
accuracy $\varepsilon/(2M)$ and confidence $1/(20M)$.
\end{definition}

\begin{definition}[Bounded Finite-Family Fourier Estimator Program]
\lean{FABL.boundedFiniteFamilyFourierEstimatorProgram}
\statementsource{lmn-constant-depth-fourier}{
  pdf-63759bd909d5-p010-b001,
  pdf-63759bd909d5-p010-b002
}
\statementsource{Analysis-of-Boolean-Functions-by-Ryan-ODonnell}{pdf-ffbe9c7b3306-p080-b002}
\label{FABL.boundedFiniteFamilyFourierEstimatorProgram}
\leanok
\uses{FABL.LearningProgram, FABL.PositiveLearningParameter, FABL.SparseFourierHypothesis, FABL.boundedFamilySampleCount, FABL.finiteFamilyFourierEstimatorProgramWithSamples}
Let $n\in\mathbb{N}$, let $\mathcal{F}$ be a finite family of subsets of the $n$
coordinates, let $0<\varepsilon\leq 1/2$ be a rational learning parameter, and let $M$
be a positive integer. Define the random-example learning program that estimates a
sparse Fourier hypothesis over $\mathcal{F}$ using the bounded-family sample count for
$\varepsilon$ and $M$.
\end{definition}

\begin{theorem}[Accuracy comparison for a bounded family]
\lean{FABL.boundedFamilyCoefficientAccuracy_le_finiteFamilyCoefficientAccuracy}
\label{FABL.boundedFamilyCoefficientAccuracy_le_finiteFamilyCoefficientAccuracy}
\leanok
\uses{FABL.PositiveLearningParameter, FABL.boundedFamilyCoefficientAccuracy, FABL.finiteFamilyCoefficientAccuracy}
\difficulty{3}
Let $\mathcal F$ be a nonempty finite family of subsets of $\{1,\ldots,n\}$, let
$0<\varepsilon\le 1/2$, and let $M$ be a positive integer with $|\mathcal F|\le M$. Then
$\varepsilon/(2M)\le\varepsilon/(2|\mathcal F|)$.
\end{theorem}

\begin{theorem}[Relative hamming dist bounded finite family output le of no bad]
\lean{FABL.relativeHammingDist_boundedFiniteFamilyOutput_le_of_no_bad}
\difficulty{3}
\statementsource{lmn-constant-depth-fourier}{
  pdf-63759bd909d5-p010-b007,
  pdf-63759bd909d5-p010-b008,
  pdf-63759bd909d5-p010-b009,
  pdf-63759bd909d5-p011-b001
}
\proofsource{Analysis-of-Boolean-Functions-by-Ryan-ODonnell}{
  pdf-ffbe9c7b3306-p080-b004,
  pdf-ffbe9c7b3306-p080-b005,
  pdf-ffbe9c7b3306-p080-b006,
  pdf-ffbe9c7b3306-p081-b001
}
\label{FABL.relativeHammingDist_boundedFiniteFamilyOutput_le_of_no_bad}
\leanok
\uses{FABL.BooleanFunction, FABL.BooleanFunction.toReal, FABL.IsFourierSpectrumConcentratedOn, FABL.PositiveLearningParameter, FABL.SparseFourierHypothesis.evaluate, FABL.SparseFourierHypothesis.ofCoefficients, FABL.boundedFamilyCoefficientAccuracy, FABL.boundedFamilyCoefficientAccuracy_le_finiteFamilyCoefficientAccuracy, FABL.boundedFamilyCoefficientConfidence, FABL.boundedFamilySampleCount, FABL.empiricalFourierCoeff, FABL.empiricalFourierCoeff_cast, FABL.finiteFamilyCoefficientBadSetWithParameters, FABL.finiteFamilyFourierEstimatorOutput, FABL.fourierCoeff, FABL.realEmpiricalFourierCoeff, FABL.relativeHammingDist, FABL.relativeHammingDist_sparseFourierHypothesis_of_coefficients_le}
% dataset-pipeline: stub -- prose pending build_informal_statements.py
Statement pending regeneration.
\end{theorem}

\begin{theorem}[Bounded finite family fourier estimator program failure probability le one twentieth]
\lean{FABL.boundedFiniteFamilyFourierEstimatorProgram_failureProbability_le_one_twentieth}
\statementsource{lmn-constant-depth-fourier}{
  pdf-63759bd909d5-p010-b004,
  pdf-63759bd909d5-p010-b005
}
\statementsource{Analysis-of-Boolean-Functions-by-Ryan-ODonnell}{
  pdf-ffbe9c7b3306-p080-b006,
  pdf-ffbe9c7b3306-p081-b001
}
\label{FABL.boundedFiniteFamilyFourierEstimatorProgram_failureProbability_le_one_twentieth}
\leanok
\uses{FABL.BooleanFunction, FABL.BooleanFunction.toReal, FABL.IsFourierSpectrumConcentratedOn, FABL.LearningProgram.eventProbability, FABL.LearningProgram.randomExampleMatrixInputLaw, FABL.LearningProgram.randomExampleMatrixOutcome, FABL.LearningProgram.randomExampleMatrixOutputLaw, FABL.LearningProgram.randomExampleMatrixOutputLaw_toOuterMeasure_apply, FABL.LearningProgram.runWithCost, FABL.PositiveLearningParameter, FABL.SparseFourierHypothesis.evaluate, FABL.boundedFamilyCoefficientAccuracy, FABL.boundedFamilyCoefficientConfidence, FABL.boundedFamilySampleCount, FABL.boundedFiniteFamilyFourierEstimatorProgram, FABL.finiteFamilyCoefficientBadSetWithParameters, FABL.finiteFamilyFourierEstimatorLabeledOutput, FABL.finiteFamilyFourierEstimatorOutput, FABL.measure_finiteFamily_someCoefficientBadSetWithParameters_le, FABL.relativeHammingDist, FABL.relativeHammingDist_boundedFiniteFamilyOutput_le_of_no_bad}
\difficulty{5}
% dataset-pipeline: stub -- prose pending build_informal_statements.py
Statement pending regeneration.
\end{theorem}

\begin{theorem}[Cost of the bounded-family Fourier estimator]
\lean{FABL.boundedFiniteFamilyFourierEstimatorProgram_cost_eq}
\label{FABL.boundedFiniteFamilyFourierEstimatorProgram_cost_eq}
\leanok
\uses{FABL.BooleanFunction, FABL.LearningCost, FABL.LearningProgram.randomExampleMatrixOutputLaw, FABL.LearningProgram.runWithCost, FABL.PositiveLearningParameter, FABL.SparseFourierHypothesis, FABL.boundedFamilySampleCount, FABL.boundedFiniteFamilyFourierEstimatorProgram, FABL.finiteFamilyFourierEstimatorProgramWithSamples}
\difficulty{4}
Let $f:\{-1,1\}^n\to\{-1,1\}$ be a sign-valued target, let $\mathcal F$ be any finite
family of subsets of $\{1,\ldots,n\}$, let $0<\varepsilon\le1/2$, and let $M$ be a
positive integer. For every outcome in the support of the estimator that uses $m$
samples per family member, where $m$ is its bounded-family sample count, the cost is
exactly $(|\mathcal F|m,\,0,\,|\mathcal F|m+|\mathcal F|m(n+1))$, recording random
examples, queries, and local work, respectively.
\end{theorem}

\begin{theorem}[Bounded finite family fourier estimator program cost le]
\lean{FABL.boundedFiniteFamilyFourierEstimatorProgram_cost_le}
\statementsource{lmn-constant-depth-fourier}{pdf-63759bd909d5-p010-b002}
\label{FABL.boundedFiniteFamilyFourierEstimatorProgram_cost_le}
\leanok
\uses{FABL.BooleanFunction, FABL.LearningCost, FABL.LearningProgram.runWithCost, FABL.PositiveLearningParameter, FABL.SparseFourierHypothesis, FABL.boundedFamilySampleCount, FABL.boundedFiniteFamilyFourierEstimatorProgram, FABL.boundedFiniteFamilyFourierEstimatorProgram_cost_eq}
\difficulty{3}
% dataset-pipeline: stub -- prose pending build_informal_statements.py
Statement pending regeneration.
\end{theorem}

\begin{definition}[Monotone Selected Fourier Family]
\lean{FABL.monotoneSelectedFourierFamily}
\statementsource{Analysis-of-Boolean-Functions-by-Ryan-ODonnell}{pdf-ffbe9c7b3306-p268-b003}
\label{FABL.monotoneSelectedFourierFamily}
\leanok
\uses{FABL.PositiveInfluenceBound, FABL.PositiveLearningParameter, FABL.learningConsequencesDegree, FABL.lowDegreeSubsets}
Let $n\in\mathbb{N}$, let $t$ be a positive rational influence bound, let
$0<\varepsilon\leq 1/2$ be a rational learning parameter, and let $J$ be a set of
coordinates among the $n$ coordinates. Define the selected Fourier family to consist of
all subsets of $J$ whose size is at most the degree prescribed by $t$ and $\varepsilon$.
\end{definition}

\begin{theorem}[Nonemptiness of the selected Fourier family]
\lean{FABL.monotoneSelectedFourierFamily_nonempty}
\label{FABL.monotoneSelectedFourierFamily_nonempty}
\leanok
\uses{FABL.PositiveInfluenceBound, FABL.PositiveLearningParameter, FABL.monotoneSelectedFourierFamily}
\difficulty{1}
Let $n$ be a nonnegative integer, let $t$ be a positive rational number, let
$0<\varepsilon\le 1/2$, and let $J$ be a subset of the $n$ coordinates. The family of
subsets of $J$ of size at most $\lceil 8t/\varepsilon\rceil$ is nonempty.
\end{theorem}

\begin{theorem}[Bound on a selected Fourier family]
\lean{FABL.card_monotoneSelectedFourierFamily_le}
\label{FABL.card_monotoneSelectedFourierFamily_le}
\leanok
\uses{FABL.PositiveInfluenceBound, FABL.PositiveLearningParameter, FABL.card_lowDegreeSubsets_le, FABL.learningConsequencesCoordinateCap, FABL.learningConsequencesDegree, FABL.learningConsequencesFamilyBound, FABL.monotoneSelectedFourierFamily}
\difficulty{3}
Let $n$ be a nonnegative integer, let $t$ be a positive rational number, and let
$\varepsilon$ be a rational number with $0<\varepsilon\le 1/2$. If $J$ is a subset of
$[n]$ whose cardinality is at most the coordinate cap determined by $t$ and
$\varepsilon$, then the family of subsets of $J$ of size at most the corresponding
degree has cardinality at most the corresponding family bound.
\end{theorem}

\begin{definition}[Monotone Bounded-Influence Second Stage]
\lean{FABL.monotoneBoundedInfluenceSecondStage}
\statementsource{Analysis-of-Boolean-Functions-by-Ryan-ODonnell}{pdf-ffbe9c7b3306-p273-b008}
\label{FABL.monotoneBoundedInfluenceSecondStage}
\leanok
\uses{FABL.LearningProgram, FABL.PositiveInfluenceBound, FABL.PositiveLearningParameter, FABL.SparseFourierHypothesis, FABL.boundedFiniteFamilyFourierEstimatorProgram, FABL.learningConsequencesFamilyBound, FABL.learningConsequencesFamilyBound_pos, FABL.monotoneSelectedFourierFamily}
Let $n\in\mathbb{N}$, let $t$ be a positive rational influence bound, let
$0<\varepsilon\leq 1/2$ be a rational learning parameter, and let $J$ be a set of
coordinates among the $n$ coordinates. Define the second-stage random-example learning
program to run the bounded finite-family Fourier estimator on the selected family of
subsets of $J$, using accuracy parameter $\varepsilon$ and the family-size bound
prescribed by $t$ and $\varepsilon$, and to return a sparse Fourier hypothesis.
\end{definition}

\begin{definition}[Monotone bounded influence learner]
\lean{FABL.monotoneBoundedInfluenceRandomExampleProgram}
\statementsource{Analysis-of-Boolean-Functions-by-Ryan-ODonnell}{
  pdf-ffbe9c7b3306-p079-b001,
  pdf-ffbe9c7b3306-p081-b005,
  pdf-ffbe9c7b3306-p081-b007,
  pdf-ffbe9c7b3306-p273-b008
}
\label{FABL.monotoneBoundedInfluenceRandomExampleProgram}
\leanok
\uses{FABL.LearningProgram, FABL.LearningProgram.bind, FABL.PositiveInfluenceBound, FABL.PositiveLearningParameter, FABL.SparseFourierHypothesis, FABL.monotoneBoundedInfluenceSecondStage, FABL.monotoneInfluenceSelectionProgram}
For each dimension $n\in\mathbb{N}$, positive rational influence bound $t$, and rational
accuracy parameter $\varepsilon$ with $0<\varepsilon\le 1/2$, this defines a
random-example learning program whose output is a sparse Fourier hypothesis on the
$n$-dimensional Boolean cube, by combining the bounded-influence second stage with the
monotone influence-selection program.
\end{definition}

\begin{theorem}[Fourier spectrum concentration]
\lean{FABL.monotoneSelectedFourierFamily_spectrum_concentrated}
\proofsource{Analysis-of-Boolean-Functions-by-Ryan-ODonnell}{
  pdf-ffbe9c7b3306-p268-b003,
  pdf-ffbe9c7b3306-p268-b004,
  pdf-ffbe9c7b3306-p269-b001
}
\label{FABL.monotoneSelectedFourierFamily_spectrum_concentrated}
\leanok
\uses{FABL.BooleanFunction, FABL.BooleanFunction.toReal, FABL.IsFourierSpectrumConcentratedOn, FABL.PositiveInfluenceBound, FABL.PositiveLearningParameter, FABL.fourierWeightOutside_antitone, FABL.influentialCoordinates, FABL.learningConsequencesDegree, FABL.learningConsequencesTargetFamily, FABL.learningConsequencesTargetFamily_spectrum_concentrated, FABL.learningConsequencesThreshold, FABL.lowDegreeSubsets, FABL.mem_lowDegreeSubsets, FABL.monotoneSelectedFourierFamily, FABL.totalInfluence}
\difficulty{4}
Let $f$ be a Boolean function on the $n$-dimensional sign cube, let $t>0$ be rational,
and let $0<\varepsilon\le 1/2$. Suppose the total influence of $f$ is at most $t$, and
let $J$ contain every coordinate whose influence is at least $\tau$, where
$\tau=(\varepsilon/8)^2t^{-2}(1/9)^{\lceil 8t/\varepsilon\rceil}$. Then the sum of the
squared Fourier coefficients of $f$ over sets not contained in the selected family of
subsets of $J$ of size at most $\lceil 8t/\varepsilon\rceil$ is at most $\varepsilon/2$.
\end{theorem}

\begin{theorem}[Failure probability for monotone targets]
\lean{FABL.monotoneBoundedInfluenceRandomExampleProgram_failureProbability_le_one_tenth}
\proofsource{Analysis-of-Boolean-Functions-by-Ryan-ODonnell}{
  pdf-ffbe9c7b3306-p273-b008,
  pdf-ffbe9c7b3306-p079-b001
}
\label{FABL.monotoneBoundedInfluenceRandomExampleProgram_failureProbability_le_one_tenth}
\leanok
\uses{FABL.BooleanFunction, FABL.BooleanFunction.toReal, FABL.LearningProgram.eventProbability, FABL.LearningProgram.eventProbability_bind_le_add, FABL.PositiveInfluenceBound, FABL.PositiveLearningParameter, FABL.SparseFourierHypothesis.evaluate, FABL.boundedFiniteFamilyFourierEstimatorProgram_failureProbability_le_one_twentieth, FABL.card_monotoneSelectedFourierFamily_le, FABL.influentialCoordinates, FABL.learningConsequencesCoordinateCap, FABL.learningConsequencesFamilyBound, FABL.learningConsequencesFamilyBound_pos, FABL.learningConsequencesThreshold, FABL.monotoneBoundedInfluenceRandomExampleProgram, FABL.monotoneBoundedInfluenceSecondStage, FABL.monotoneInfluenceSelectionProgram, FABL.monotoneInfluenceSelectionProgram_failureProbability_le_one_twentieth, FABL.monotoneSelectedFourierFamily, FABL.monotoneSelectedFourierFamily_nonempty, FABL.monotoneSelectedFourierFamily_spectrum_concentrated, FABL.monotoneSelectionBad, FABL.relativeHammingDist, FABL.totalInfluence}
\difficulty{5}
Let $n$ be a nonnegative integer, let $f:\{\pm1\}^n\to\{\pm1\}$ be monotone, let $t>0$
and $0<\varepsilon\le\tfrac12$ be rational, and suppose the total influence of $f$ is at
most $t$. When the program is run on target $f$, the probability that its output
hypothesis has relative Hamming distance strictly greater than $\varepsilon$ from $f$ is
at most $1/10$.
\end{theorem}

\begin{definition}[Random example sample budget]
\lean{FABL.monotoneBoundedInfluenceRandomExampleBudget}
\statementsource{Analysis-of-Boolean-Functions-by-Ryan-ODonnell}{
  pdf-ffbe9c7b3306-p081-b007,
  pdf-ffbe9c7b3306-p273-b008
}
\label{FABL.monotoneBoundedInfluenceRandomExampleBudget}
\leanok
\uses{FABL.PositiveInfluenceBound, FABL.PositiveLearningParameter, FABL.boundedFamilySampleCount, FABL.learningConsequencesFamilyBound, FABL.learningConsequencesFamilyBound_pos, FABL.singletonInfluenceSampleCount}
For each $n\in\mathbb{N}$, positive rational influence bound $t$, and rational accuracy
parameter $\varepsilon$ with $0<\varepsilon\le 1/2$, define the sample budget to be $n$
times the singleton-influence sample count, plus the learning-consequences family bound
times the bounded-family sample count evaluated at that family bound and its positivity
proof.
\end{definition}

\begin{definition}[Monotone bounded influence work budget]
\lean{FABL.monotoneBoundedInfluenceWorkBudget}
\statementsource{Analysis-of-Boolean-Functions-by-Ryan-ODonnell}{
  pdf-ffbe9c7b3306-p081-b007,
  pdf-ffbe9c7b3306-p273-b008
}
\label{FABL.monotoneBoundedInfluenceWorkBudget}
\leanok
\uses{FABL.PositiveInfluenceBound, FABL.PositiveLearningParameter, FABL.boundedFamilySampleCount, FABL.learningConsequencesFamilyBound, FABL.learningConsequencesFamilyBound_pos, FABL.singletonInfluenceSampleCount}
For each $n\in\mathbb{N}$, positive rational influence bound $t$, and rational accuracy
parameter $\varepsilon$ with $0<\varepsilon\le 1/2$, define the work budget as $n$ times
the singleton-influence sample count times $n+2$, plus $n$, plus the
learning-consequences family bound times the bounded-family sample count at that bound,
times $n+2$.
\end{definition}

\begin{theorem}[Cost bounds for the learning program]
\lean{FABL.monotoneBoundedInfluenceRandomExampleProgram_cost_le}
\label{FABL.monotoneBoundedInfluenceRandomExampleProgram_cost_le}
\leanok
\uses{FABL.BooleanFunction, FABL.LearningCost, FABL.LearningProgram.randomExampleMatrixOutputLaw, FABL.LearningProgram.runWithCost, FABL.LearningProgram.runWithCost_bind, FABL.PositiveInfluenceBound, FABL.PositiveLearningParameter, FABL.SparseFourierHypothesis, FABL.boundedFamilySampleCount, FABL.boundedFiniteFamilyFourierEstimatorProgram_cost_le, FABL.card_empiricalInfluentialCoordinates_le, FABL.card_monotoneSelectedFourierFamily_le, FABL.learningConsequencesCoordinateCap, FABL.learningConsequencesFamilyBound, FABL.learningConsequencesFamilyBound_pos, FABL.monotoneBoundedInfluenceRandomExampleBudget, FABL.monotoneBoundedInfluenceRandomExampleProgram, FABL.monotoneBoundedInfluenceWorkBudget, FABL.monotoneInfluenceSelectionProgram, FABL.monotoneInfluenceSelectionProgram_cost_eq, FABL.monotoneSelectedFourierFamily, FABL.singletonInfluenceSampleCount}
\difficulty{5}
Let $n$ be a nonnegative integer, let $t>0$ and $0<\varepsilon\le\tfrac12$ be rational,
and let $f:\{\pm1\}^n\to\{\pm1\}$ be a Boolean function. Every outcome in the support of
the program's cost distribution on target $f$ uses at most the monotone
bounded-influence random-example budget, makes exactly zero membership queries, and
incurs work at most the monotone bounded-influence work budget, with parameters
$n,t,\varepsilon$.
\end{theorem}

\begin{theorem}[Leaf count under output mapping]
\lean{FABL.FiniteDecisionTree.leafCount_mapOutput}
\label{FABL.FiniteDecisionTree.leafCount_mapOutput}
\leanok
\uses{FABL.FiniteDecisionTree, FABL.FiniteDecisionTree.leafCount, FABL.FiniteDecisionTree.mapOutput}
\difficulty{4}
Let $T$ be a finite decision tree with queries from a finite alphabet, and let $g$ be
any map from its output labels to another set of labels. The tree obtained by applying
$g$ to every leaf label has the same number of leaves as $T$.
\end{theorem}

\begin{theorem}[Output relabeling preserves computation]
\lean{FABL.FiniteDecisionTree.mapOutput_signValue_computes}
\label{FABL.FiniteDecisionTree.mapOutput_signValue_computes}
\leanok
\uses{FABL.BooleanFunction, FABL.BooleanFunction.toReal, FABL.FiniteAlphabetDecisionTree, FABL.FiniteDecisionTree.Computes, FABL.FiniteDecisionTree.eval_mapOutput, FABL.FiniteDecisionTree.mapOutput, FABL.Sign, FABL.signValue}
\difficulty{3}
Let $n$ be a natural number, let $f : \{-1,1\}^n \to \{\pm1\}$ be a Boolean
function, and let $T$ be a finite decision tree that queries coordinates of inputs in
$\{\pm1\}^n$ and outputs signs. If $T$ computes $f$, then the tree obtained by
replacing each leaf sign with its corresponding real value computes the real-valued
function obtained by viewing each value of $f$ as a real number.
\end{theorem}

\begin{theorem}[Total influence le sqrt logb leaf count of monotone]
\lean{FABL.totalInfluence_le_sqrt_logb_leafCount_of_monotone}
\statementsource{Analysis-of-Boolean-Functions-by-Ryan-ODonnell}{
  pdf-ffbe9c7b3306-p232-b005,
  pdf-ffbe9c7b3306-p232-b006,
  pdf-ffbe9c7b3306-p232-b008,
  pdf-ffbe9c7b3306-p243-b002
}
\label{FABL.totalInfluence_le_sqrt_logb_leafCount_of_monotone}
\leanok
\uses{FABL.BooleanFunction, FABL.BooleanFunction.toReal, FABL.FiniteAlphabetDecisionTree, FABL.FiniteDecisionTree.Computes, FABL.FiniteDecisionTree.deterministicAverageCost_le_log_leafCount, FABL.FiniteDecisionTree.leafCount, FABL.FiniteDecisionTree.leafCount_mapOutput, FABL.FiniteDecisionTree.mapOutput, FABL.FiniteDecisionTree.mapOutput_signValue_computes, FABL.RandomizedDecisionTree.averageCost, FABL.RandomizedDecisionTree.averageCost_ofDeterministic, FABL.RandomizedDecisionTree.deterministicAverageCost, FABL.RandomizedDecisionTree.ofDeterministic, FABL.RandomizedDecisionTree.totalInfluence_le_sigma_mul_sqrt_averageCost, FABL.Sign, FABL.biasSigma, FABL.biasVarianceScale, FABL.biasedSignPMF, FABL.biasedSignPMF_one_half_eq_uniform, FABL.productTotalInfluence, FABL.productTotalInfluence_uniformSign_eq_totalInfluence, FABL.signValue, FABL.totalInfluence, FABL.uniformPMF}
\difficulty{5}
% dataset-pipeline: stub -- prose pending build_informal_statements.py
Statement pending regeneration.
\end{theorem}

\begin{theorem}[Tree attaining deterministic size]
\lean{FABL.FiniteDecisionTree.exists_computingTree_leafCount_eq_deterministicSize}
\label{FABL.FiniteDecisionTree.exists_computingTree_leafCount_eq_deterministicSize}
\leanok
\uses{FABL.BooleanFunction, FABL.FiniteAlphabetDecisionTree, FABL.FiniteDecisionTree.Computes, FABL.FiniteDecisionTree.deterministicSize, FABL.FiniteDecisionTree.exists_computingTree_with_leafCount, FABL.FiniteDecisionTree.leafCount, FABL.Sign}
\difficulty{1}
Let $n$ be a natural number and let $f:\{−1,1\}^n\to\{−1,1\}$ be a Boolean function.
There exists a decision tree that computes $f$ and whose number of leaves equals the
least number of leaves of any decision tree computing $f$.
\end{theorem}

\begin{theorem}[Influence bound for monotone functions]
\lean{FABL.totalInfluence_le_sqrt_logb_deterministicSize_of_monotone}
\label{FABL.totalInfluence_le_sqrt_logb_deterministicSize_of_monotone}
\leanok
\uses{FABL.BooleanFunction, FABL.BooleanFunction.toReal, FABL.FiniteDecisionTree.deterministicSize, FABL.FiniteDecisionTree.exists_computingTree_leafCount_eq_deterministicSize, FABL.totalInfluence, FABL.totalInfluence_le_sqrt_logb_leafCount_of_monotone}
\difficulty{4}
Let $n$ be a natural number and let $f$ be a monotone Boolean function on the
$n$-dimensional sign cube. Its total influence is at most $\sqrt{\log_2 s}$, where $s$
is the minimum number of leaves in a deterministic decision tree computing $f$.
\end{theorem}

\begin{definition}[Decision tree size influence bound]
\lean{FABL.decisionTreeSizeInfluenceBound}
\statementsource{Analysis-of-Boolean-Functions-by-Ryan-ODonnell}{
  pdf-ffbe9c7b3306-p243-b002,
  pdf-ffbe9c7b3306-p273-b008
}
\label{FABL.decisionTreeSizeInfluenceBound}
\leanok
\uses{FABL.PositiveInfluenceBound}
For natural numbers $n$ and $d$, define the positive rational influence bound
$t=\sqrt{d\,\lceil\log_2(n+1)\rceil}+1$, where the square root is rounded down to an
integer.
\end{definition}

\begin{theorem}[Total influence le decision tree size influence bound]
\lean{FABL.totalInfluence_le_decisionTreeSizeInfluenceBound}
\statementsource{Analysis-of-Boolean-Functions-by-Ryan-ODonnell}{
  pdf-ffbe9c7b3306-p232-b005,
  pdf-ffbe9c7b3306-p232-b006,
  pdf-ffbe9c7b3306-p232-b008,
  pdf-ffbe9c7b3306-p243-b002
}
\label{FABL.totalInfluence_le_decisionTreeSizeInfluenceBound}
\leanok
\uses{FABL.BooleanFunction, FABL.BooleanFunction.toReal, FABL.FiniteDecisionTree.deterministicSize, FABL.FiniteDecisionTree.exists_computingTree_leafCount_eq_deterministicSize, FABL.FiniteDecisionTree.leafCount_pos', FABL.decisionTreeSizeInfluenceBound, FABL.totalInfluence, FABL.totalInfluence_le_sqrt_logb_deterministicSize_of_monotone}
\difficulty{5}
% dataset-pipeline: stub -- prose pending build_informal_statements.py
Statement pending regeneration.
\end{theorem}

\begin{definition}[Monotone decision-tree learning program]
\lean{FABL.monotoneDecisionTreeSizeRandomExampleProgram}
\statementsource{Analysis-of-Boolean-Functions-by-Ryan-ODonnell}{
  pdf-ffbe9c7b3306-p079-b002,
  pdf-ffbe9c7b3306-p082-b004,
  pdf-ffbe9c7b3306-p273-b008
}
\label{FABL.monotoneDecisionTreeSizeRandomExampleProgram}
\leanok
\uses{FABL.LearningProgram, FABL.PositiveLearningParameter, FABL.SparseFourierHypothesis, FABL.decisionTreeSizeInfluenceBound, FABL.monotoneBoundedInfluenceRandomExampleProgram}
For natural numbers $n$ and $d$, and a rational parameter $0<\varepsilon\le 1/2$, define
a random-example learning program that returns a sparse Fourier hypothesis on the
$n$-dimensional Boolean cube. It is the bounded-influence random-example program with
the influence bound associated with decision-tree size at most $d$ and parameter
$\varepsilon$.
\end{definition}

\begin{theorem}[Learning failure probability]
\lean{FABL.monotoneDecisionTreeSizeRandomExampleProgram_failureProbability_le_one_tenth}
\label{FABL.monotoneDecisionTreeSizeRandomExampleProgram_failureProbability_le_one_tenth}
\leanok
\uses{FABL.BooleanFunction, FABL.FiniteDecisionTree.deterministicSize, FABL.LearningProgram.eventProbability, FABL.PositiveLearningParameter, FABL.SparseFourierHypothesis.evaluate, FABL.decisionTreeSizeInfluenceBound, FABL.monotoneBoundedInfluenceRandomExampleProgram_failureProbability_le_one_tenth, FABL.monotoneDecisionTreeSizeRandomExampleProgram, FABL.relativeHammingDist, FABL.totalInfluence_le_decisionTreeSizeInfluenceBound}
\difficulty{1}
Let $n$ be a nonnegative integer, let $f$ be a monotone Boolean function on the
$n$-dimensional sign cube, and suppose its deterministic decision-tree size is at most
$n^d$ for some nonnegative integer $d$. For every $0<\varepsilon\le 1/2$, the
probability that the program using the decision-tree-size influence bound outputs a
hypothesis whose relative Hamming distance from $f$ is greater than $\varepsilon$ is at
most $1/10$.
\end{theorem}

\begin{theorem}[Monotone decision tree size random example program cost le]
\lean{FABL.monotoneDecisionTreeSizeRandomExampleProgram_cost_le}
\statementsource{Analysis-of-Boolean-Functions-by-Ryan-ODonnell}{
  pdf-ffbe9c7b3306-p082-b004,
  pdf-ffbe9c7b3306-p273-b008
}
\label{FABL.monotoneDecisionTreeSizeRandomExampleProgram_cost_le}
\leanok
\uses{FABL.BooleanFunction, FABL.LearningCost, FABL.LearningProgram.runWithCost, FABL.PositiveLearningParameter, FABL.SparseFourierHypothesis, FABL.decisionTreeSizeInfluenceBound, FABL.monotoneBoundedInfluenceRandomExampleBudget, FABL.monotoneBoundedInfluenceRandomExampleProgram_cost_le, FABL.monotoneBoundedInfluenceWorkBudget, FABL.monotoneDecisionTreeSizeRandomExampleProgram}
\difficulty{3}
% dataset-pipeline: stub -- prose pending build_informal_statements.py
Statement pending regeneration.
\end{theorem}
